\par\bigskip\pdfbookmark[0]{On the Connection. Fragment V}{nexus-fragment-V}\centerline{[V.]}\medskip
The average number of classes\begingroup\renewcommand{\thefootnote}{*)}\footnote{[Compare \emph{Disquiss. Arithm.}, Article 302; the formula there differs by a constant $\delta$ from the one occurring here in the text.]}\endgroup\ near the negative determinant $-D$ is approximately
\[
=\frac{\pi\surd D}{4(1+\tfrac1{27}+\tfrac1{125}+\text{etc.})}
\]
The actual number is expressed by the following formulas, where, for brevity, $m$ stands for the average number and $M$ for the actual number; $p,q$ denote all odd primes that do not divide $D$, the former being divisors and the latter nondivisors of $\square+D$; $r$ denotes the primes\begingroup\renewcommand{\thefootnote}{**)}\footnote{[odd.]}\endgroup\ that divide $D$:

\SourceLeaf{295}
\[
\begin{aligned}
\text{I.}\quad M&=m\ \text{Product of }\frac{p^3+p^2}{p^3-1}\pdot\frac{q^3-q^2}{q^3-1}\pdot\frac{r^3-r}{r^3-1}\\
\text{II.}\quad M&=\frac{\pi\surd D}4\ \text{Product of }\frac{p+1}p\pdot\frac{q-1}q\pdot\frac{rr-1}{rr}\\
\text{III. NB.}\quad M&=\frac{2\surd D}\pi\ \text{Product of }\frac p{p-1}\pdot\frac q{q+1}\\
\text{IV.}\quad M&=\surd\left\{\frac D2\pdot\text{Product of }\frac{p+1}{p-1}\pdot\frac{q-1}{q+1}\pdot\frac{rr-1}{rr}\right\}\\
\text{V.}\quad M&=\frac{2\surd D}\pi\{1\pm\tfrac13\pm\tfrac15\pm\tfrac17\pm\tfrac19\text{ etc.}\}
\end{aligned}
\]
NB. Formula III follows directly by comparing the two ways of counting the representable numbers up to a given bound.
\BookRule{.30\textwidth}
